% Complete candidate replacements after independent reading of all four drafts.
% Numeric checks: check_delivered.py and check_fable.py in this directory.
% Historical feasibility witnesses and semantic labels remain inherited.

\section*{Proposed abstract}

When an agent sells to a competitor, the proceeds improve its own finances
while the asset may strengthen its rival. We study how four configured LLM
systems negotiate this trade-off in twenty complete four-player Monopoly
games. A census of 920 offer responses connects bargaining to subsequent
construction, debt payment and inheritance. Agents explicitly consider rival
development funds, intermediary resale and prospective estates. Among opening
offers, recipients reject 138 of 187 exchanges completing only the proposer,
compared with 152 of 238 completing neither side. Their own resources also
matter: holders below \$100 reject 17 of 69 completion requests, versus
120 of 175 at \$300 or more; low-cash holders mostly counter rather than
accept. Thirteen of fourteen completing exchanges that follow an earlier
decline by the same holder occur after its cash has fallen. Immediate
financing is usually executed successfully: the systems pay 128 of 131 debts
with demonstrated payment capacity, including sixty-one multistep payments.
Eight contextual incidents connect these patterns to stated strategies and
observed actions. Provider refusals alter Opus's note and retry conditions,
and some Grok calls conflict with accompanying text. The results describe
these configured systems in this cohort. They reveal resource-dependent
bargaining and concrete financial accomplishments without identifying
optimal policies or ranking underlying models.

\section*{Scientific motivation and contribution}

A completing deed has two economic roles. It is an asset for its holder and
an opportunity for the player who owns the rest of its group. Selling it can
fund the holder's construction while enabling rent against that same holder.
An intermediary can also buy the deed to resell it, and a potential creditor
may prefer to inherit it rather than prolong its owner's survival. Negotiation
in this setting requires attention to consequences beyond the immediate
objects and cash exchanged.

We examine that attention in twenty complete games played by four configured
LLM systems. The scientific object is their behavior: what considerations
appear in offers and explanations, how bargaining varies with current
resources, and which financial tasks the systems actually complete. The
records connect each offer and response to delivered state and later events,
allowing us to follow economic positions across negotiation, development and
bankruptcy rather than treating a deal or a win as a sufficient measure of
strategy.

Three observations organize the study. First, agents bargain over the
opportunities transactions create for rivals, including immediate construction,
third-party resale and future estates. Second, holders with little cash are
more likely to continue bargaining over completing deeds, usually by
counteroffering. Third, immediate debt finance is often successful even though
the development and income of acquired groups differ sharply across the
observed trajectories. These are linked behavioral findings: the holder's
resources, the rival's opportunity and the execution of the resulting position
all enter the same economic interaction.

\section*{Opening offers and continued concessions}

Of the 918 offers that could legally be accepted, 289 would complete only the
author's group, 140 only the recipient's, 81 both parties' groups and 408
neither. Thread position changes the response distribution. Among opening
offers, recipients reject 138 of 187 author-only completions and 152 of 238
non-completing exchanges. The 9.9 percentage-point gap has the same direction
in all five seed blocks and for Sol, Gemini and Grok, while Opus reverses.
The comparison concerns different transactions and circumstances; it describes
the observed market rather than isolating a completion effect.

Continued offers provide another view. We define a cash concession as an
author improving the recipient's net cash while preserving the noncash terms
of its own previous offer in that thread. Recipients accept 9 of 58 such
offers completing only the proposer and 24 of 86 completing neither side.
These observations occur in 38 and 43 threads respectively. The lower
acceptance of rival-completing concessions shows that bargaining can persist
despite improved cash terms, with the dependence among offers retained in
the interpretation.

The explanations make some of these stakes explicit. Both earlier
LLM-assisted readings identify a trade-linked rival consequence in 61 of
74 refusals of the recipient's own completion, including 27 of 38 cases in
which the rival would gain no new group. The concerns include development
funds, retained blockers and other holdings. These agreed codes are
concentrated: Gemini and Grok supply 56 of the 61, and G1 and G15 supply 31.
They document expressed considerations, with no human-validated measure of
whether the resulting decisions were correct.

\section*{Cash pressure changes the form of bargaining}

For requests that would complete only the proposer, holders with less than
\$100 reject 17 of 69, counter 44 and accept 8. Holders with at least \$300
reject 120 of 175, counter 48 and accept 7. Thus the main change is from
outright refusal to continued negotiation. The rejection ordering persists
within opening and later offers and in each of the five seed blocks. It is
not uniform across configurations: Gemini rejects ten of its eleven low-cash
requests, while Opus rarely rejects at either cash level.

We then follow the same potential buyer, holder and color group within a
game. Fourteen exchanges completing only one party follow a refusal or
counter by that holder on an earlier turn. In thirteen, the holder has less
cash at settlement; median cash is \$571.50 at the first decline and \$82
at settlement. The transferred completing deeds are the same, but reciprocal
deeds change in nine of the fourteen cases, and cash terms and board positions
also change. These selected trajectories show how earlier resistance can
give way to agreement as circumstances evolve. They do not establish a
threshold at which a holder must sell.

Cash pressure and sale price are separate. In G7, Gemini renews a bid after
Sol spends on dark-blue development. Sol's note relates its counterprice to
the houses it can finance and the buyer's remaining funds. The resulting
\$700 sale completes Gemini's oranges; Gemini builds four orange houses
in the same turn and Sol buys three more dark-blue houses within eight turns.
Selling a rival its completing deed can therefore coexist with a stated
development plan of the seller's own. The observation is the plan, exchange
and subsequent actions, not an estimated benefit of that sale over keeping
the deed.

\section*{Prices and the limits of a valuation claim}

Seller asks are higher for completing buyers in many seller, mortgage and
thread-position strata. The closer overlaps are scarce and mixed. Seven
game--seller--deed groups contain asks to both completing and non-completing
buyers: median asks are higher in four, lower in two and equal in one.
Using the first ask changes the split to five, one and one because G4's
completing asks fall during negotiation. Buyer bids differ by configuration.
The evidence supports attention to a buyer's position without identifying
a general completion premium or a calibrated valuation rule.

This distinction matters for the economic interpretation. Monopoly's list
price provides a common scale, not the deed's value to a particular player.
An exchange below list is not by itself a fire sale, and one above list does
not disprove financial pressure. Similarly, a nominal post-trade building
budget omits alternative collateral, redemption, reserves and group quality.
The present comparisons do not determine how accurately agents value those
options.

\section*{Positive evidence of financial execution}

The cohort contains 188 episodes in which a debt exceeds cash and the player
receives a liquidation menu. Of these, 128 end in payment and 60 in
elimination. Sixty-one successful payments require two to ten actions.
Across the paid episodes, the applied sequences contain 180 mortgage actions
and 89 building-sale actions, with no engine fallback. In G20, Sol combines
\$49 cash, \$700 in building salvage and \$665 in mortgages to meet a
\$1,400 debt and retain \$14.

Three eliminated positions have previously verified legal payment sequences.
Together with the observed payments, they give 131 episodes with demonstrated
payment capacity, of which 128 are paid. Fifty-six other eliminations exceed
a generous own-asset bound, and the remaining G7 position is restricted by
its actual liquidation menu. This census supplies positive evidence about
executing immediate financing through the provided interface. It does not
require treating a note as access to internal computation or a later win as
proof that earlier investments were good.

\section*{Development and economic consequences}

Trade-completed groups become income only when their owner can build, chooses
to finance the build and submits an action the system applies. Sol and Gemini
build in the first management turn on 33 of 35 completed groups; Opus and
Grok do so on 10 of 28. The latter trajectories include response-resolution
conditions as well as deferred redemption and stated holds. Group ownership
durations and aggregate opponent-landing rates are similar across the
first-window classes, while realized rent per landing differs substantially.
This makes the development history necessary context for the trade count.
It does not hold group quality, landing timing or financial position fixed.

The seller's subsequent position also matters. Among 41 one-sided completing
exchanges that deliver net cash to the holder, thirteen are followed by that
holder's construction within eight turns. Nine involve an eventual winner,
compared with one among the other twenty-eight. These are exchange-level
trajectories selected by later construction; one winning player appears in
both categories. They motivate attention to what follows a sale rather than
establish a causal advantage of selling or building.

\section*{Proposed conclusion}

The observed agents bargain over rival opportunities as well as their own
receipts. They discuss construction funds, intermediary resale and prospective
estates, and their responses change with the holder's current resources.
Low-cash holders usually reopen negotiation through counters rather than
immediately release a completing deed. The systems also accomplish a concrete
financial task with high observed success: 128 of 131 demonstrated payable
debt episodes end in payment. Development and inheritance connect these local
decisions to the rest of the game. Taken together, the records show strategic
attention, resource-dependent bargaining and successful financial execution
under the studied conditions, while leaving optimal play and model rankings
unresolved.
